% Part A. Included from methods.tex with % Part A. Included from methods.tex with % Part A. Included from methods.tex with % Part A. Included from methods.tex with \input{sections/partA_methods}.
\subsection{Preprocessing (Part A)}
Part A works on the \SI{1}{\kilo\hertz} means of each channel and has three steps:
\begin{enumerate}
  \item a median filter (3 samples), which removes single-sample spikes;
  \item a Savitzky--Golay filter (7 samples, order 3)~\cite{savitzky1964}, which smooths small
        fluctuations;
  \item subtraction of a slowly varying baseline. The baseline is fitted by asymmetric least
        squares (ALS)~\cite{eilers2005} on \SI{10}{\hertz} block means (smoothness
        $\lambda = 10^{5}$, asymmetry $p = 0.01$, 10 iterations) and interpolated back to
        \SI{1}{\kilo\hertz}. Samples above the fitted curve get a small weight, so the baseline
        follows the signal from below and is not pulled up by the peaks.
\end{enumerate}
The ALS equations form a symmetric pentadiagonal system, which is solved with a banded
Cholesky solver, and the two channels are processed in parallel threads. On all recordings
compared, this gives the same pieces and labels as the earlier sparse-matrix implementation.
If the median ratio \texttt{adc1}/\texttt{adc2} at the 50 largest \texttt{adc2} peaks is below
0.1, \texttt{adc1} is not used for the two-channel check.

\subsection{Noise level and threshold}
\label{sec:threshold}
The noise level $\sigma$ of each channel is estimated from the cleaned signal itself. The signal
is divided into \SI{200}{\milli\second} blocks, blocks whose standard deviation is below half the
median block standard deviation are dropped as idle (no flow), and $\sigma$ is the root mean
square of the standard deviations of the quietest 10\,\% of the remaining blocks.
Part B detects peaks whose prominence exceeds $k\sigma$ with $k = 10$.
All settings are the same for every recording: $k = 10$, and the grouping gap and the no-dye
level given with Part B. Only $\sigma$ differs between recordings, and it is measured from each
signal and not chosen by concentration.
.
\subsection{Preprocessing (Part A)}
Part A works on the \SI{1}{\kilo\hertz} means of each channel and has three steps:
\begin{enumerate}
  \item a median filter (3 samples), which removes single-sample spikes;
  \item a Savitzky--Golay filter (7 samples, order 3)~\cite{savitzky1964}, which smooths small
        fluctuations;
  \item subtraction of a slowly varying baseline. The baseline is fitted by asymmetric least
        squares (ALS)~\cite{eilers2005} on \SI{10}{\hertz} block means (smoothness
        $\lambda = 10^{5}$, asymmetry $p = 0.01$, 10 iterations) and interpolated back to
        \SI{1}{\kilo\hertz}. Samples above the fitted curve get a small weight, so the baseline
        follows the signal from below and is not pulled up by the peaks.
\end{enumerate}
The ALS equations form a symmetric pentadiagonal system, which is solved with a banded
Cholesky solver, and the two channels are processed in parallel threads. On all recordings
compared, this gives the same pieces and labels as the earlier sparse-matrix implementation.
If the median ratio \texttt{adc1}/\texttt{adc2} at the 50 largest \texttt{adc2} peaks is below
0.1, \texttt{adc1} is not used for the two-channel check.

\subsection{Noise level and threshold}
\label{sec:threshold}
The noise level $\sigma$ of each channel is estimated from the cleaned signal itself. The signal
is divided into \SI{200}{\milli\second} blocks, blocks whose standard deviation is below half the
median block standard deviation are dropped as idle (no flow), and $\sigma$ is the root mean
square of the standard deviations of the quietest 10\,\% of the remaining blocks.
Part B detects peaks whose prominence exceeds $k\sigma$ with $k = 10$.
All settings are the same for every recording: $k = 10$, and the grouping gap and the no-dye
level given with Part B. Only $\sigma$ differs between recordings, and it is measured from each
signal and not chosen by concentration.
.
\subsection{Preprocessing (Part A)}
Part A works on the \SI{1}{\kilo\hertz} means of each channel and has three steps:
\begin{enumerate}
  \item a median filter (3 samples), which removes single-sample spikes;
  \item a Savitzky--Golay filter (7 samples, order 3)~\cite{savitzky1964}, which smooths small
        fluctuations;
  \item subtraction of a slowly varying baseline. The baseline is fitted by asymmetric least
        squares (ALS)~\cite{eilers2005} on \SI{10}{\hertz} block means (smoothness
        $\lambda = 10^{5}$, asymmetry $p = 0.01$, 10 iterations) and interpolated back to
        \SI{1}{\kilo\hertz}. Samples above the fitted curve get a small weight, so the baseline
        follows the signal from below and is not pulled up by the peaks.
\end{enumerate}
The ALS equations form a symmetric pentadiagonal system, which is solved with a banded
Cholesky solver, and the two channels are processed in parallel threads. On all recordings
compared, this gives the same pieces and labels as the earlier sparse-matrix implementation.
If the median ratio \texttt{adc1}/\texttt{adc2} at the 50 largest \texttt{adc2} peaks is below
0.1, \texttt{adc1} is not used for the two-channel check.

\subsection{Noise level and threshold}
\label{sec:threshold}
The noise level $\sigma$ of each channel is estimated from the cleaned signal itself. The signal
is divided into \SI{200}{\milli\second} blocks, blocks whose standard deviation is below half the
median block standard deviation are dropped as idle (no flow), and $\sigma$ is the root mean
square of the standard deviations of the quietest 10\,\% of the remaining blocks.
Part B detects peaks whose prominence exceeds $k\sigma$ with $k = 10$.
All settings are the same for every recording: $k = 10$, and the grouping gap and the no-dye
level given with Part B. Only $\sigma$ differs between recordings, and it is measured from each
signal and not chosen by concentration.
.
\subsection{Preprocessing (Part A)}
Part A works on the \SI{1}{\kilo\hertz} means of each channel and has three steps:
\begin{enumerate}
  \item a median filter (3 samples), which removes single-sample spikes;
  \item a Savitzky--Golay filter (7 samples, order 3)~\cite{savitzky1964}, which smooths small
        fluctuations;
  \item subtraction of a slowly varying baseline. The baseline is fitted by asymmetric least
        squares (ALS)~\cite{eilers2005} on \SI{10}{\hertz} block means (smoothness
        $\lambda = 10^{5}$, asymmetry $p = 0.01$, 10 iterations) and interpolated back to
        \SI{1}{\kilo\hertz}. Samples above the fitted curve get a small weight, so the baseline
        follows the signal from below and is not pulled up by the peaks.
\end{enumerate}
The ALS equations form a symmetric pentadiagonal system, which is solved with a banded
Cholesky solver, and the two channels are processed in parallel threads. On all recordings
compared, this gives the same pieces and labels as the earlier sparse-matrix implementation.
If the median ratio \texttt{adc1}/\texttt{adc2} at the 50 largest \texttt{adc2} peaks is below
0.1, \texttt{adc1} is not used for the two-channel check.

\subsection{Noise level and threshold}
\label{sec:threshold}
The noise level $\sigma$ of each channel is estimated from the cleaned signal itself. The signal
is divided into \SI{200}{\milli\second} blocks, blocks whose standard deviation is below half the
median block standard deviation are dropped as idle (no flow), and $\sigma$ is the root mean
square of the standard deviations of the quietest 10\,\% of the remaining blocks.
Part B detects peaks whose prominence exceeds $k\sigma$ with $k = 10$.
All settings are the same for every recording: $k = 10$, and the grouping gap and the no-dye
level given with Part B. Only $\sigma$ differs between recordings, and it is measured from each
signal and not chosen by concentration.
